% BHU PYQ -- Manually transcribed & error-corrected
% Paper: BCF-122 : Fundamentals of Business Finance
% Exam: B.Com. (Hons.) FMM Semester II Examination, 2021-22

\documentclass[11pt,a4paper]{article}
\usepackage[a4paper,top=2.5cm,bottom=2.5cm,left=2.8cm,right=2.8cm,headheight=14pt]{geometry}
\usepackage{amsmath,amssymb}
\usepackage[T1]{fontenc}
\usepackage[utf8]{inputenc}
\usepackage{lmodern,microtype,fancyhdr,enumitem,hyperref,lastpage,parskip}

\hypersetup{
    colorlinks=true,
    urlcolor=black,
    linkcolor=black,
    pdftitle={BCF-122: Fundamentals of Business Finance -- BHU B.Com. (Hons.) FMM Sem II 2021-22}
}

\pagestyle{fancy}
\fancyhf{}
\renewcommand{\headrulewidth}{0.4pt}
\renewcommand{\footrulewidth}{0.4pt}
\fancyhead[L]{\small   \textit{Banaras Hindu University}}
\fancyhead[R]{\small   \textit{BCF-122: Fundamentals of Business Finance}}
\fancyfoot[C]{\small Page  \thepage\ of \pageref{LastPage}}


\newlist{parts}{enumerate}{2}
\setlist[parts,1]{label=(\alph*),leftmargin=*,itemsep=5pt,topsep=3pt}
\setlist[parts,2]{label=(\roman*),leftmargin=*,itemsep=3pt,topsep=2pt}

\newcommand{\pts}[1]{\hfill   \textit{[#1]}}


\begin{document}


\begin{center}
  {\large\bfseries BANARAS HINDU UNIVERSITY}\\[2pt]
  {\normalsize B.Com. (Hons.) FMM --- Semester II Examination, 2021-22}\\[6pt]
  \rule{0.8\linewidth}{0.5pt}\\[6pt]
  {\large\bfseries Paper No. BCF-122: Fundamentals of Business Finance}\\[4pt]
  {\normalsize Time: 3 Hours\hfill Maximum Marks: 70}
\end{center}

\rule{\linewidth}{0.4pt}

\smallskip



\noindent   \textit{Instructions: Answer all questions. All questions carry equal marks (14 marks each).}

\bigskip




\noindent   \textbf{Question 1.}
Describe the concept of business finance and explain its nature and scope.\pts{14}

\centerline{   \textbf{OR}}

Explain the various limitations of the profit maximization concept. How can we overcome these through the wealth maximization concept?\pts{14}

\medskip\rule{\linewidth}{0.2pt}




\noindent   \textbf{Question 2.}
What is financial planning? Explain its nature and types.\pts{14}

\centerline{   \textbf{OR}}

What do you understand by `financial decision-making'? Throw light on its various aspects.\pts{14}

\medskip\rule{\linewidth}{0.2pt}




\noindent   \textbf{Question 3.}
What are the determinants of the capital structure of a company? Does an optimum capital structure exist in an organization?\pts{14}

\centerline{   \textbf{OR}}

Draw points of difference between over-capitalization and under-capitalization in an organization.\pts{14}

\medskip\rule{\linewidth}{0.2pt}




\noindent   \textbf{Question 4.}
Define the concept of cost of capital. How would you determine the weighted average cost of capital of a firm?\pts{14}

\centerline{   \textbf{OR}}

Write short notes on any two of the following:\pts{14}

\begin{parts}
  \item Cost of equity
  \item Cost of debt capital
  \item Cost of retained earnings
\end{parts}

\medskip\rule{\linewidth}{0.2pt}




\noindent   \textbf{Question 5.}
What do you understand by capital budgeting? Discuss its process and importance.\pts{14}

\centerline{   \textbf{OR}}

What is a dividend? Describe the factors affecting dividend policies.\pts{14}

\vfill
\rule{\linewidth}{0.4pt}

\begin{center}\small
\href{https://bhu-exam.vercel.app/}{Click here to visit our website}\\[4pt]
\href{https://me-aryan.vercel.app/}{Click here to visit our contributor}\\[4pt]
\href{https://quashan-me.vercel.app/}{Click here to visit our contributor}
\end{center}

\end{document}
